
\documentclass{article}
\usepackage{colortbl}
\usepackage{makecell}
\usepackage{multirow}
\usepackage{supertabular}

\begin{document}

\newcounter{utterance}

\centering \large Interaction Transcript for game `clean\_up', experiment `0\_easy\_7obj\_lt', episode 1 with claude{-}opus{-}4{-}8{-}azure.
\vspace{24pt}

{ \footnotesize  \setcounter{utterance}{1}
\setlength{\tabcolsep}{0pt}
\begin{supertabular}{c@{$\;$}|p{.15\linewidth}@{}p{.15\linewidth}p{.15\linewidth}p{.15\linewidth}p{.15\linewidth}p{.15\linewidth}}
    \# & $\;$A & \multicolumn{4}{c}{Game Master} & $\;\:$B\\
    \hline

    \theutterance \stepcounter{utterance}  
    & & \multicolumn{4}{p{0.6\linewidth}}{
        \cellcolor[rgb]{0.9,0.9,0.9}{
            \makecell[{{p{\linewidth}}}]{
                \texttt{\tiny{[P1$\langle$GM]}}
                \texttt{Aš esu jūsų žaidimo vedėjas, o jūs žaidžiate bendrą žaidimą, kurio žaidimo lenta yra tokia lentelė:} \\
\\ 
\texttt{```} \\
\texttt{ 1234567} \\
\texttt{╔═══════╗} \\
\texttt{║◌◌◌◌◌T◌║ 1} \\
\texttt{║P◌◌◌◌◌◌║ 2} \\
\texttt{║◌◌AD◌┌─╢ 3} \\
\texttt{║◌◌◌┌─┤S║ 4} \\
\texttt{║O◌◌│◌├─╢ 5} \\
\texttt{╟───┼─┘◌║ 6} \\
\texttt{║◌M◌│◌◌◌║ 7} \\
\texttt{╚═══╧═══╝} \\
\\ 
\texttt{```} \\
\\ 
\texttt{* Viršutiniame krašte x koordinatės didėja į dešinę, o dešiniajame krašte y koordinatės didėja žemyn.} \\
\texttt{* Tinklelyje atsitiktine tvarka išdėstyti šie objektai: 'P', 'O', 'T', 'S', 'D', 'A', 'M'.} \\
\\ 
\\ 
\texttt{Kitas žaidėjas mato kitą žaidimo lentos variantą, kuriame objektai išdėstyti skirtingose atsitiktinėse vietose. Jūs negalite matyti kito žaidėjo lentos, o jis negali matyti jūsų lentos.} \\
\\ 
\\ 
\texttt{**Tikslas:**} \\
\\ 
\texttt{Abu žaidėjai turi perkelti objektus savo fone taip, kad vienodi objektai atsidurtų tose pačiose koordinatėse. Turite bendrauti su kitu žaidėju ir susitarti dėl bendros tikslinės konfigūracijos.} \\
\\ 
\\ 
\texttt{**Taisyklės:**} \\
\\ 
\texttt{* Kiekvieno ėjimo metu galite išsiųsti lygiai vieną iš šių dviejų komandų:} \\
\texttt{ 1. `SAKYTI: <ŽINUTĖ>`: norėdami išsiųsti žinutę (viską iki kitos eilutės pertraukos) kitam žaidėjui. Aš ją persiųsiu jūsų partneriui.} \\
\texttt{2. `JUDĖTI: <OBJEKTAS>, (<X>, <Y>)`: perkelti objektą į naują poziciją, kur `<X>` yra stulpelis, o `<Y>` yra eilutė. Aš jums pranešiu, ar jūsų perkėlimas buvo sėkmingas, ar ne.} \\
\texttt{* Jei nesilaikysite formato arba siųsite kelias komandas iš karto, turėsiu jus nubausti.} \\
\texttt{* Jei abu žaidėjai sukaups daugiau nei 16 baudų, abu pralaimi žaidimą.} \\
\texttt{* Labai svarbu, kad bendrautumėte su kitu žaidėju dėl savo tikslinės būsenos! Vienintelis būdas, kuriuo galite perduoti savo strategiją kitam žaidėjui, yra komanda `SAKYTI: <ŽINUTĖ>`!} \\
\\ 
\\ 
\texttt{**Objektų perkėlimas**} \\
\\ 
\texttt{* Objektus galima perkelti tik į langelius, esančius tinklelio ribose. Tikslinis langelis turi būti tuščias, t. y. jame turi būti tik simbolis '◌'.} \\
\texttt{* Jei bandysite perkelti objektą į vietą, kuri nėra tuščia, arba bandysite jį perkelti už tinklelio ribų, turėsiu jus nubausti. Gausite dar vieną bandymą.} \\
\texttt{* Prieš atlikdami perkėlimą, dukart patikrinkite, ar tikslinė vieta yra tuščia ir ar joje nėra jokios raidės, rėmelio ar linijos!} \\
\\ 
\\ 
\texttt{**Žaidimo pabaiga**} \\
\\ 
\texttt{Jei manote, kad pasiekėte tikslą {-} sulygiavote visus objektus, galite paprašyti kito žaidėjo baigti žaidimą siųsdami `SAKYTI: baigta?`. Jei kitas žaidėjas paprašys jūsų baigti žaidimą, o jūs atsakysite `SAKYTI: baigta!`, žaidimas baigsis.} \\
\\ 
\texttt{Abu žaidėjai laimi, jei žaidimas baigiasi per 28 raundus, kur vienas raundas apibrėžiamas kaip du žaidėjai, kiekvienas išsiuntęs po vieną galiojančią komandą.} \\
\\ 
\texttt{**Taškų skaičiavimas:**} \\
\\ 
\texttt{Kuo arčiau vienodų objektų abiejose žaidimo lentose, tuo daugiau taškų gausite. Baudos mažina jūsų taškus. Ar pavyks pagerinti rekordą?} \\
\\ 
\\ 
\texttt{Išsiųskite žinutę kitam žaidėjui, kad pradėtumėte žaidimą!} \\
            }
        }
    }
    & \\ \\

    \theutterance \stepcounter{utterance}  
    & \multicolumn{4}{p{0.6\linewidth}}{
        \cellcolor[rgb]{0.9,0.9,0.9}{
            \makecell[{{p{\linewidth}}}]{
                \texttt{\tiny{[P1$\rangle$GM]}}
                \texttt{SAKYTI: Sveiki! Siūlau bendrą tikslą: sudėkime visus 7 objektus į viršutinę eilutę (y=1) abėcėlės tvarka. Taigi: A→(1,1), D→(2,1), M→(3,1), O→(4,1), P→(5,1), S→(6,1), T→(7,1). Ar tau tinka? Ar visi šie langeliai tavo lentoje tušti?} \\
            }
        }
    }
    & & \\ \\

    \theutterance \stepcounter{utterance}  
    & & & \multicolumn{2}{p{0.3\linewidth}}{
        \cellcolor[rgb]{0.9,0.9,0.9}{
            \makecell[{{p{\linewidth}}}]{
                \texttt{\tiny{[GM$|$GM]}}
                \texttt{SAKYTI: Sveiki! Siūlau bendrą tikslą: sudėkime visus 7 objektus į viršutinę eilutę (y=1) abėcėlės tvarka. Taigi: A→(1,1), D→(2,1), M→(3,1), O→(4,1), P→(5,1), S→(6,1), T→(7,1). Ar tau tinka? Ar visi šie langeliai tavo lentoje tušti?} \\
            }
        }
    }
    & & \\ \\

    \theutterance \stepcounter{utterance}  
    & & \multicolumn{4}{p{0.6\linewidth}}{
        \cellcolor[rgb]{0.9,0.9,0.9}{
            \makecell[{{p{\linewidth}}}]{
                \texttt{\tiny{[GM$\rangle$P2]}}
                \texttt{Aš esu jūsų žaidimo vedėjas, o jūs žaidžiate bendrą žaidimą, kurio žaidimo lenta yra tokia lentelė:} \\
\\ 
\texttt{```} \\
\texttt{ 1234567} \\
\texttt{╔═══════╗} \\
\texttt{║◌◌◌◌◌S◌║ 1} \\
\texttt{║◌◌◌◌O◌◌║ 2} \\
\texttt{║◌◌◌M◌┌─╢ 3} \\
\texttt{║◌◌A┌─┤◌║ 4} \\
\texttt{║◌◌◌│◌├─╢ 5} \\
\texttt{╟───┼─┘◌║ 6} \\
\texttt{║D◌T│P◌◌║ 7} \\
\texttt{╚═══╧═══╝} \\
\\ 
\texttt{```} \\
\\ 
\texttt{* Viršutiniame krašte x koordinatės didėja į dešinę, o dešiniajame krašte y koordinatės didėja žemyn.} \\
\texttt{* Tinklelyje atsitiktine tvarka išdėstyti šie objektai: 'P', 'O', 'T', 'S', 'D', 'A', 'M'.} \\
\\ 
\\ 
\texttt{Kitas žaidėjas mato kitą žaidimo lentos variantą, kuriame objektai išdėstyti skirtingose atsitiktinėse vietose. Jūs negalite matyti kito žaidėjo lentos, o jis negali matyti jūsų lentos.} \\
\\ 
\\ 
\texttt{**Tikslas:**} \\
\\ 
\texttt{Abu žaidėjai turi perkelti objektus savo fone taip, kad vienodi objektai atsidurtų tose pačiose koordinatėse. Turite bendrauti su kitu žaidėju ir susitarti dėl bendros tikslinės konfigūracijos.} \\
\\ 
\\ 
\texttt{**Taisyklės:**} \\
\\ 
\texttt{* Kiekvieno ėjimo metu galite išsiųsti lygiai vieną iš šių dviejų komandų:} \\
\texttt{ 1. `SAKYTI: <ŽINUTĖ>`: norėdami išsiųsti žinutę (viską iki kitos eilutės pertraukos) kitam žaidėjui. Aš ją persiųsiu jūsų partneriui.} \\
\texttt{2. `JUDĖTI: <OBJEKTAS>, (<X>, <Y>)`: perkelti objektą į naują poziciją, kur `<X>` yra stulpelis, o `<Y>` yra eilutė. Aš jums pranešiu, ar jūsų perkėlimas buvo sėkmingas, ar ne.} \\
\texttt{* Jei nesilaikysite formato arba siųsite kelias komandas iš karto, turėsiu jus nubausti.} \\
\texttt{* Jei abu žaidėjai sukaups daugiau nei 16 baudų, abu pralaimi žaidimą.} \\
\texttt{* Labai svarbu, kad bendrautumėte su kitu žaidėju dėl savo tikslinės būsenos! Vienintelis būdas, kuriuo galite perduoti savo strategiją kitam žaidėjui, yra komanda `SAKYTI: <ŽINUTĖ>`!} \\
\\ 
\\ 
\texttt{**Objektų perkėlimas**} \\
\\ 
\texttt{* Objektus galima perkelti tik į langelius, esančius tinklelio ribose. Tikslinis langelis turi būti tuščias, t. y. jame turi būti tik simbolis '◌'.} \\
\texttt{* Jei bandysite perkelti objektą į vietą, kuri nėra tuščia, arba bandysite jį perkelti už tinklelio ribų, turėsiu jus nubausti. Gausite dar vieną bandymą.} \\
\texttt{* Prieš atlikdami perkėlimą, dukart patikrinkite, ar tikslinė vieta yra tuščia ir ar joje nėra jokios raidės, rėmelio ar linijos!} \\
\\ 
\\ 
\texttt{**Žaidimo pabaiga**} \\
\\ 
\texttt{Jei manote, kad pasiekėte tikslą {-} sulygiavote visus objektus, galite paprašyti kito žaidėjo baigti žaidimą siųsdami `SAKYTI: baigta?`. Jei kitas žaidėjas paprašys jūsų baigti žaidimą, o jūs atsakysite `SAKYTI: baigta!`, žaidimas baigsis.} \\
\\ 
\texttt{Abu žaidėjai laimi, jei žaidimas baigiasi per 28 raundus, kur vienas raundas apibrėžiamas kaip du žaidėjai, kiekvienas išsiuntęs po vieną galiojančią komandą.} \\
\\ 
\texttt{**Taškų skaičiavimas:**} \\
\\ 
\texttt{Kuo arčiau vienodų objektų abiejose žaidimo lentose, tuo daugiau taškų gausite. Baudos mažina jūsų taškus. Ar pavyks pagerinti rekordą?} \\
\\ 
\texttt{Kitas žaidėjas pradėjo žaidimą išsiųsdamas šią žinutę:} \\
\\ 
\texttt{"Sveiki! Siūlau bendrą tikslą: sudėkime visus 7 objektus į viršutinę eilutę (y=1) abėcėlės tvarka. Taigi: A→(1,1), D→(2,1), M→(3,1), O→(4,1), P→(5,1), S→(6,1), T→(7,1). Ar tau tinka? Ar visi šie langeliai tavo lentoje tušti?"} \\
\\ 
\texttt{Kokia yra jūsų pirmoji komanda?} \\
            }
        }
    }
    & \\ \\

    \theutterance \stepcounter{utterance}  
    & & & \multicolumn{4}{p{0.6\linewidth}}{
        \cellcolor[rgb]{0.9,0.9,0.9}{
            \makecell[{{p{\linewidth}}}]{
                \texttt{\tiny{[GM$\langle$P2]}}
                \texttt{Pažiūrėkime į mano lentą ir viršutinę eilutę (y=1):} \\
\\ 
\texttt{{-} (1,1)=◌, (2,1)=◌, (3,1)=◌, (4,1)=◌, (5,1)=◌, (6,1)=**S**, (7,1)=◌} \\
\\ 
\texttt{Visi viršutinės eilutės langeliai tušti, o **S** jau stovi (6,1) — kaip tik ten, kur reikia pagal pasiūlymą. Tai puikiai tinka.} \\
\\ 
\texttt{Mano objektų dabartinės pozicijos:} \\
\texttt{{-} A (3,4), D (1,7), M (4,3), O (5,2), P (5,7), S (6,1), T (3,7)} \\
\\ 
\texttt{Sutinku su pasiūlymu, tad pirmiausia patvirtinsiu:} \\
\\ 
\texttt{**SAKYTI: Tinka! Visi viršutinės eilutės langeliai mano lentoje tušti, o S jau stovi (6,1). Pradedu judinti kitus objektus į vietas: A→(1,1), D→(2,1), M→(3,1), O→(4,1), P→(5,1), T→(7,1).**} \\
            }
        }
    }
     \\ \\

    \theutterance \stepcounter{utterance}  
    & & & \multicolumn{2}{p{0.3\linewidth}}{
        \cellcolor[rgb]{0.9,0.9,0.9}{
            \makecell[{{p{\linewidth}}}]{
                \texttt{\tiny{[GM$|$GM]}}
                \texttt{Pažiūrėkime į mano lentą ir viršutinę eilutę (y=1):} \\
\\ 
\texttt{{-} (1,1)=◌, (2,1)=◌, (3,1)=◌, (4,1)=◌, (5,1)=◌, (6,1)=**S**, (7,1)=◌} \\
\\ 
\texttt{Visi viršutinės eilutės langeliai tušti, o **S** jau stovi (6,1) — kaip tik ten, kur reikia pagal pasiūlymą. Tai puikiai tinka.} \\
\\ 
\texttt{Mano objektų dabartinės pozicijos:} \\
\texttt{{-} A (3,4), D (1,7), M (4,3), O (5,2), P (5,7), S (6,1), T (3,7)} \\
\\ 
\texttt{Sutinku su pasiūlymu, tad pirmiausia patvirtinsiu:} \\
\\ 
\texttt{**SAKYTI: Tinka! Visi viršutinės eilutės langeliai mano lentoje tušti, o S jau stovi (6,1). Pradedu judinti kitus objektus į vietas: A→(1,1), D→(2,1), M→(3,1), O→(4,1), P→(5,1), T→(7,1).**} \\
            }
        }
    }
    & & \\ \\

    \theutterance \stepcounter{utterance}  
    & & & \multicolumn{2}{p{0.3\linewidth}}{
        \cellcolor[rgb]{0.9,0.9,0.9}{
            \makecell[{{p{\linewidth}}}]{
                \texttt{\tiny{[GM$|$GM]}}
                \texttt{Invalid format: head is not empty: 'Pažiūrėkime į mano lentą ir viršutinę eilutę (y=1):} \\
\\ 
\texttt{{-} (1,1)=◌, (2,1)=◌, (3,1)=◌, (4,1)=◌, (5,1)=◌, (6,1)=**S**, (7,1)=◌} \\
\\ 
\texttt{Visi viršutinės eilutės langeliai tušti, o **S** jau stovi (6,1) — kaip tik ten, kur reikia pagal pasiūlymą. Tai puikiai tinka.} \\
\\ 
\texttt{Mano objektų dabartinės pozicijos:} \\
\texttt{{-} A (3,4), D (1,7), M (4,3), O (5,2), P (5,7), S (6,1), T (3,7)} \\
\\ 
\texttt{Sutinku su pasiūlymu, tad pirmiausia patvirtinsiu:} \\
\\ 
\texttt{**'} \\
            }
        }
    }
    & & \\ \\

    \theutterance \stepcounter{utterance}  
    & & \multicolumn{4}{p{0.6\linewidth}}{
        \cellcolor[rgb]{0.9,0.9,0.9}{
            \makecell[{{p{\linewidth}}}]{
                \texttt{\tiny{[GM$\rangle$P2]}}
                \texttt{Nuobauda: Jūsų žinutėje neturi būti nieko prieš komandą!} \\
\\ 
\texttt{Įsitikinkite, kad jūsų atsakyme yra tik `SAKYTI: <ŽINUTĖ>` arba `JUDĖTI: <OBJEKTAS>, (<X>, <Y>)`, ir nieko daugiau!} \\
\\ 
\texttt{Jūs bendrai sukaupėte 1 iš 16 baudų.} \\
\texttt{Pabandykite dar kartą!} \\
            }
        }
    }
    & \\ \\

    \theutterance \stepcounter{utterance}  
    & & & \multicolumn{4}{p{0.6\linewidth}}{
        \cellcolor[rgb]{0.9,0.9,0.9}{
            \makecell[{{p{\linewidth}}}]{
                \texttt{\tiny{[GM$\langle$P2]}}
                \texttt{SAKYTI: Tinka! Visi viršutinės eilutės langeliai mano lentoje tušti, o S jau stovi (6,1). Pradedu judinti kitus objektus į vietas.} \\
            }
        }
    }
     \\ \\

    \theutterance \stepcounter{utterance}  
    & & & \multicolumn{2}{p{0.3\linewidth}}{
        \cellcolor[rgb]{0.9,0.9,0.9}{
            \makecell[{{p{\linewidth}}}]{
                \texttt{\tiny{[GM$|$GM]}}
                \texttt{SAKYTI: Tinka! Visi viršutinės eilutės langeliai mano lentoje tušti, o S jau stovi (6,1). Pradedu judinti kitus objektus į vietas.} \\
            }
        }
    }
    & & \\ \\

    \theutterance \stepcounter{utterance}  
    & & \multicolumn{4}{p{0.6\linewidth}}{
        \cellcolor[rgb]{0.9,0.9,0.9}{
            \makecell[{{p{\linewidth}}}]{
                \texttt{\tiny{[P1$\langle$GM]}}
                \texttt{Jūsų žinutė buvo perduota kitam žaidėjui.} \\
\\ 
\\ 
\texttt{Šiuo metu žaidžiate 1 raundą iš maksimaliai 28.} \\
\texttt{Jūs bendrai sukaupėte 1 iš 16 baudų.} \\
\texttt{Kitas žaidėjas atsiuntė šią žinutę:} \\
\\ 
\\ 
\texttt{"Tinka! Visi viršutinės eilutės langeliai mano lentoje tušti, o S jau stovi (6,1). Pradedu judinti kitus objektus į vietas."} \\
\\ 
\\ 
\texttt{Kokia jūsų kita komanda?} \\
            }
        }
    }
    & \\ \\

    \theutterance \stepcounter{utterance}  
    & \multicolumn{4}{p{0.6\linewidth}}{
        \cellcolor[rgb]{0.9,0.9,0.9}{
            \makecell[{{p{\linewidth}}}]{
                \texttt{\tiny{[P1$\rangle$GM]}}
                \texttt{SAKYTI: Puiku! Pradedu judinti. Pirmiausia perkelsiu T į (7,1), kad atlaisvinčiau (6,1) dėl S.} \\
            }
        }
    }
    & & \\ \\

    \theutterance \stepcounter{utterance}  
    & & & \multicolumn{2}{p{0.3\linewidth}}{
        \cellcolor[rgb]{0.9,0.9,0.9}{
            \makecell[{{p{\linewidth}}}]{
                \texttt{\tiny{[GM$|$GM]}}
                \texttt{SAKYTI: Puiku! Pradedu judinti. Pirmiausia perkelsiu T į (7,1), kad atlaisvinčiau (6,1) dėl S.} \\
            }
        }
    }
    & & \\ \\

    \theutterance \stepcounter{utterance}  
    & & \multicolumn{4}{p{0.6\linewidth}}{
        \cellcolor[rgb]{0.9,0.9,0.9}{
            \makecell[{{p{\linewidth}}}]{
                \texttt{\tiny{[GM$\rangle$P2]}}
                \texttt{Jūsų žinutė buvo perduota kitam žaidėjui.} \\
\\ 
\\ 
\texttt{Šiuo metu žaidžiate 2 raundą iš maksimaliai 28.} \\
\texttt{Jūs bendrai sukaupėte 1 iš 16 baudų.} \\
\texttt{Kitas žaidėjas atsiuntė šią žinutę:} \\
\\ 
\\ 
\texttt{"Puiku! Pradedu judinti. Pirmiausia perkelsiu T į (7,1), kad atlaisvinčiau (6,1) dėl S."} \\
\\ 
\\ 
\texttt{Kokia jūsų kita komanda?} \\
            }
        }
    }
    & \\ \\

    \theutterance \stepcounter{utterance}  
    & & & \multicolumn{4}{p{0.6\linewidth}}{
        \cellcolor[rgb]{0.9,0.9,0.9}{
            \makecell[{{p{\linewidth}}}]{
                \texttt{\tiny{[GM$\langle$P2]}}
                \texttt{JUDĖTI: A, (1, 1)} \\
            }
        }
    }
     \\ \\

    \theutterance \stepcounter{utterance}  
    & & & \multicolumn{2}{p{0.3\linewidth}}{
        \cellcolor[rgb]{0.9,0.9,0.9}{
            \makecell[{{p{\linewidth}}}]{
                \texttt{\tiny{[GM$|$GM]}}
                \texttt{JUDĖTI: A, (1, 1)} \\
            }
        }
    }
    & & \\ \\

    \theutterance \stepcounter{utterance}  
    & & & \multicolumn{2}{p{0.3\linewidth}}{
        \cellcolor[rgb]{0.9,0.9,0.9}{
            \makecell[{{p{\linewidth}}}]{
                \texttt{\tiny{[GM$|$GM]}}
                \texttt{A moved to (1, 1)} \\
            }
        }
    }
    & & \\ \\

    \theutterance \stepcounter{utterance}  
    & & \multicolumn{4}{p{0.6\linewidth}}{
        \cellcolor[rgb]{0.9,0.9,0.9}{
            \makecell[{{p{\linewidth}}}]{
                \texttt{\tiny{[P1$\langle$GM]}}
                \texttt{Jūsų žinutė buvo perduota kitam žaidėjui.} \\
\\ 
\\ 
\texttt{Šiuo metu žaidžiate 2 raundą iš maksimaliai 28.} \\
\texttt{Jūs bendrai sukaupėte 1 iš 16 baudų.} \\
\texttt{Kitas žaidėjas perkėlė objektą savo tinklelyje.} \\
\\ 
\\ 
\texttt{Kokia jūsų kita komanda?} \\
            }
        }
    }
    & \\ \\

    \theutterance \stepcounter{utterance}  
    & \multicolumn{4}{p{0.6\linewidth}}{
        \cellcolor[rgb]{0.9,0.9,0.9}{
            \makecell[{{p{\linewidth}}}]{
                \texttt{\tiny{[P1$\rangle$GM]}}
                \texttt{JUDĖTI: T, (7, 1)} \\
            }
        }
    }
    & & \\ \\

    \theutterance \stepcounter{utterance}  
    & & & \multicolumn{2}{p{0.3\linewidth}}{
        \cellcolor[rgb]{0.9,0.9,0.9}{
            \makecell[{{p{\linewidth}}}]{
                \texttt{\tiny{[GM$|$GM]}}
                \texttt{JUDĖTI: T, (7, 1)} \\
            }
        }
    }
    & & \\ \\

    \theutterance \stepcounter{utterance}  
    & & & \multicolumn{2}{p{0.3\linewidth}}{
        \cellcolor[rgb]{0.9,0.9,0.9}{
            \makecell[{{p{\linewidth}}}]{
                \texttt{\tiny{[GM$|$GM]}}
                \texttt{T moved to (7, 1)} \\
            }
        }
    }
    & & \\ \\

    \theutterance \stepcounter{utterance}  
    & & \multicolumn{4}{p{0.6\linewidth}}{
        \cellcolor[rgb]{0.9,0.9,0.9}{
            \makecell[{{p{\linewidth}}}]{
                \texttt{\tiny{[GM$\rangle$P2]}}
                \texttt{Sėkmingai perkelta 'A' į (1,1). Jūsų atnaujintas tinklelis atrodo taip:} \\
\texttt{```} \\
\texttt{ 1234567} \\
\texttt{╔═══════╗} \\
\texttt{║A◌◌◌◌S◌║ 1} \\
\texttt{║◌◌◌◌O◌◌║ 2} \\
\texttt{║◌◌◌M◌┌─╢ 3} \\
\texttt{║◌◌◌┌─┤◌║ 4} \\
\texttt{║◌◌◌│◌├─╢ 5} \\
\texttt{╟───┼─┘◌║ 6} \\
\texttt{║D◌T│P◌◌║ 7} \\
\texttt{╚═══╧═══╝} \\
\texttt{```} \\
\texttt{Šiuo metu žaidžiate 3 raundą iš maksimaliai 28.} \\
\texttt{Jūs bendrai sukaupėte 1 iš 16 baudų.} \\
\texttt{Kitas žaidėjas perkėlė objektą savo tinklelyje.} \\
\\ 
\\ 
\texttt{Kokia jūsų kita komanda?} \\
            }
        }
    }
    & \\ \\

    \theutterance \stepcounter{utterance}  
    & & & \multicolumn{4}{p{0.6\linewidth}}{
        \cellcolor[rgb]{0.9,0.9,0.9}{
            \makecell[{{p{\linewidth}}}]{
                \texttt{\tiny{[GM$\langle$P2]}}
                \texttt{JUDĖTI: D, (2, 1)} \\
            }
        }
    }
     \\ \\

    \theutterance \stepcounter{utterance}  
    & & & \multicolumn{2}{p{0.3\linewidth}}{
        \cellcolor[rgb]{0.9,0.9,0.9}{
            \makecell[{{p{\linewidth}}}]{
                \texttt{\tiny{[GM$|$GM]}}
                \texttt{JUDĖTI: D, (2, 1)} \\
            }
        }
    }
    & & \\ \\

    \theutterance \stepcounter{utterance}  
    & & & \multicolumn{2}{p{0.3\linewidth}}{
        \cellcolor[rgb]{0.9,0.9,0.9}{
            \makecell[{{p{\linewidth}}}]{
                \texttt{\tiny{[GM$|$GM]}}
                \texttt{D moved to (2, 1)} \\
            }
        }
    }
    & & \\ \\

    \theutterance \stepcounter{utterance}  
    & & \multicolumn{4}{p{0.6\linewidth}}{
        \cellcolor[rgb]{0.9,0.9,0.9}{
            \makecell[{{p{\linewidth}}}]{
                \texttt{\tiny{[P1$\langle$GM]}}
                \texttt{Sėkmingai perkelta 'T' į (7,1). Jūsų atnaujintas tinklelis atrodo taip:} \\
\texttt{```} \\
\texttt{ 1234567} \\
\texttt{╔═══════╗} \\
\texttt{║◌◌◌◌◌◌T║ 1} \\
\texttt{║P◌◌◌◌◌◌║ 2} \\
\texttt{║◌◌AD◌┌─╢ 3} \\
\texttt{║◌◌◌┌─┤S║ 4} \\
\texttt{║O◌◌│◌├─╢ 5} \\
\texttt{╟───┼─┘◌║ 6} \\
\texttt{║◌M◌│◌◌◌║ 7} \\
\texttt{╚═══╧═══╝} \\
\texttt{```} \\
\texttt{Šiuo metu žaidžiate 3 raundą iš maksimaliai 28.} \\
\texttt{Jūs bendrai sukaupėte 1 iš 16 baudų.} \\
\texttt{Kitas žaidėjas perkėlė objektą savo tinklelyje.} \\
\\ 
\\ 
\texttt{Kokia jūsų kita komanda?} \\
            }
        }
    }
    & \\ \\

    \theutterance \stepcounter{utterance}  
    & \multicolumn{4}{p{0.6\linewidth}}{
        \cellcolor[rgb]{0.9,0.9,0.9}{
            \makecell[{{p{\linewidth}}}]{
                \texttt{\tiny{[P1$\rangle$GM]}}
                \texttt{JUDĖTI: S, (6, 1)} \\
            }
        }
    }
    & & \\ \\

    \theutterance \stepcounter{utterance}  
    & & & \multicolumn{2}{p{0.3\linewidth}}{
        \cellcolor[rgb]{0.9,0.9,0.9}{
            \makecell[{{p{\linewidth}}}]{
                \texttt{\tiny{[GM$|$GM]}}
                \texttt{JUDĖTI: S, (6, 1)} \\
            }
        }
    }
    & & \\ \\

    \theutterance \stepcounter{utterance}  
    & & & \multicolumn{2}{p{0.3\linewidth}}{
        \cellcolor[rgb]{0.9,0.9,0.9}{
            \makecell[{{p{\linewidth}}}]{
                \texttt{\tiny{[GM$|$GM]}}
                \texttt{S moved to (6, 1)} \\
            }
        }
    }
    & & \\ \\

    \theutterance \stepcounter{utterance}  
    & & \multicolumn{4}{p{0.6\linewidth}}{
        \cellcolor[rgb]{0.9,0.9,0.9}{
            \makecell[{{p{\linewidth}}}]{
                \texttt{\tiny{[GM$\rangle$P2]}}
                \texttt{Sėkmingai perkelta 'D' į (2,1). Jūsų atnaujintas tinklelis atrodo taip:} \\
\texttt{```} \\
\texttt{ 1234567} \\
\texttt{╔═══════╗} \\
\texttt{║AD◌◌◌S◌║ 1} \\
\texttt{║◌◌◌◌O◌◌║ 2} \\
\texttt{║◌◌◌M◌┌─╢ 3} \\
\texttt{║◌◌◌┌─┤◌║ 4} \\
\texttt{║◌◌◌│◌├─╢ 5} \\
\texttt{╟───┼─┘◌║ 6} \\
\texttt{║◌◌T│P◌◌║ 7} \\
\texttt{╚═══╧═══╝} \\
\texttt{```} \\
\texttt{Šiuo metu žaidžiate 4 raundą iš maksimaliai 28.} \\
\texttt{Jūs bendrai sukaupėte 1 iš 16 baudų.} \\
\texttt{Kitas žaidėjas perkėlė objektą savo tinklelyje.} \\
\\ 
\\ 
\texttt{Kokia jūsų kita komanda?} \\
            }
        }
    }
    & \\ \\

    \theutterance \stepcounter{utterance}  
    & & & \multicolumn{4}{p{0.6\linewidth}}{
        \cellcolor[rgb]{0.9,0.9,0.9}{
            \makecell[{{p{\linewidth}}}]{
                \texttt{\tiny{[GM$\langle$P2]}}
                \texttt{JUDĖTI: M, (3, 1)} \\
            }
        }
    }
     \\ \\

    \theutterance \stepcounter{utterance}  
    & & & \multicolumn{2}{p{0.3\linewidth}}{
        \cellcolor[rgb]{0.9,0.9,0.9}{
            \makecell[{{p{\linewidth}}}]{
                \texttt{\tiny{[GM$|$GM]}}
                \texttt{JUDĖTI: M, (3, 1)} \\
            }
        }
    }
    & & \\ \\

    \theutterance \stepcounter{utterance}  
    & & & \multicolumn{2}{p{0.3\linewidth}}{
        \cellcolor[rgb]{0.9,0.9,0.9}{
            \makecell[{{p{\linewidth}}}]{
                \texttt{\tiny{[GM$|$GM]}}
                \texttt{M moved to (3, 1)} \\
            }
        }
    }
    & & \\ \\

    \theutterance \stepcounter{utterance}  
    & & \multicolumn{4}{p{0.6\linewidth}}{
        \cellcolor[rgb]{0.9,0.9,0.9}{
            \makecell[{{p{\linewidth}}}]{
                \texttt{\tiny{[P1$\langle$GM]}}
                \texttt{Sėkmingai perkelta 'S' į (6,1). Jūsų atnaujintas tinklelis atrodo taip:} \\
\texttt{```} \\
\texttt{ 1234567} \\
\texttt{╔═══════╗} \\
\texttt{║◌◌◌◌◌ST║ 1} \\
\texttt{║P◌◌◌◌◌◌║ 2} \\
\texttt{║◌◌AD◌┌─╢ 3} \\
\texttt{║◌◌◌┌─┤◌║ 4} \\
\texttt{║O◌◌│◌├─╢ 5} \\
\texttt{╟───┼─┘◌║ 6} \\
\texttt{║◌M◌│◌◌◌║ 7} \\
\texttt{╚═══╧═══╝} \\
\texttt{```} \\
\texttt{Šiuo metu žaidžiate 4 raundą iš maksimaliai 28.} \\
\texttt{Jūs bendrai sukaupėte 1 iš 16 baudų.} \\
\texttt{Kitas žaidėjas perkėlė objektą savo tinklelyje.} \\
\\ 
\\ 
\texttt{Kokia jūsų kita komanda?} \\
            }
        }
    }
    & \\ \\

    \theutterance \stepcounter{utterance}  
    & \multicolumn{4}{p{0.6\linewidth}}{
        \cellcolor[rgb]{0.9,0.9,0.9}{
            \makecell[{{p{\linewidth}}}]{
                \texttt{\tiny{[P1$\rangle$GM]}}
                \texttt{JUDĖTI: A, (1, 1)} \\
            }
        }
    }
    & & \\ \\

    \theutterance \stepcounter{utterance}  
    & & & \multicolumn{2}{p{0.3\linewidth}}{
        \cellcolor[rgb]{0.9,0.9,0.9}{
            \makecell[{{p{\linewidth}}}]{
                \texttt{\tiny{[GM$|$GM]}}
                \texttt{JUDĖTI: A, (1, 1)} \\
            }
        }
    }
    & & \\ \\

    \theutterance \stepcounter{utterance}  
    & & & \multicolumn{2}{p{0.3\linewidth}}{
        \cellcolor[rgb]{0.9,0.9,0.9}{
            \makecell[{{p{\linewidth}}}]{
                \texttt{\tiny{[GM$|$GM]}}
                \texttt{A moved to (1, 1)} \\
            }
        }
    }
    & & \\ \\

    \theutterance \stepcounter{utterance}  
    & & \multicolumn{4}{p{0.6\linewidth}}{
        \cellcolor[rgb]{0.9,0.9,0.9}{
            \makecell[{{p{\linewidth}}}]{
                \texttt{\tiny{[GM$\rangle$P2]}}
                \texttt{Sėkmingai perkelta 'M' į (3,1). Jūsų atnaujintas tinklelis atrodo taip:} \\
\texttt{```} \\
\texttt{ 1234567} \\
\texttt{╔═══════╗} \\
\texttt{║ADM◌◌S◌║ 1} \\
\texttt{║◌◌◌◌O◌◌║ 2} \\
\texttt{║◌◌◌◌◌┌─╢ 3} \\
\texttt{║◌◌◌┌─┤◌║ 4} \\
\texttt{║◌◌◌│◌├─╢ 5} \\
\texttt{╟───┼─┘◌║ 6} \\
\texttt{║◌◌T│P◌◌║ 7} \\
\texttt{╚═══╧═══╝} \\
\texttt{```} \\
\texttt{Šiuo metu žaidžiate 5 raundą iš maksimaliai 28.} \\
\texttt{Jūs bendrai sukaupėte 1 iš 16 baudų.} \\
\texttt{Kitas žaidėjas perkėlė objektą savo tinklelyje.} \\
\\ 
\\ 
\texttt{Kokia jūsų kita komanda?} \\
            }
        }
    }
    & \\ \\

    \theutterance \stepcounter{utterance}  
    & & & \multicolumn{4}{p{0.6\linewidth}}{
        \cellcolor[rgb]{0.9,0.9,0.9}{
            \makecell[{{p{\linewidth}}}]{
                \texttt{\tiny{[GM$\langle$P2]}}
                \texttt{JUDĖTI: O, (4, 1)} \\
            }
        }
    }
     \\ \\

    \theutterance \stepcounter{utterance}  
    & & & \multicolumn{2}{p{0.3\linewidth}}{
        \cellcolor[rgb]{0.9,0.9,0.9}{
            \makecell[{{p{\linewidth}}}]{
                \texttt{\tiny{[GM$|$GM]}}
                \texttt{JUDĖTI: O, (4, 1)} \\
            }
        }
    }
    & & \\ \\

    \theutterance \stepcounter{utterance}  
    & & & \multicolumn{2}{p{0.3\linewidth}}{
        \cellcolor[rgb]{0.9,0.9,0.9}{
            \makecell[{{p{\linewidth}}}]{
                \texttt{\tiny{[GM$|$GM]}}
                \texttt{O moved to (4, 1)} \\
            }
        }
    }
    & & \\ \\

    \theutterance \stepcounter{utterance}  
    & & \multicolumn{4}{p{0.6\linewidth}}{
        \cellcolor[rgb]{0.9,0.9,0.9}{
            \makecell[{{p{\linewidth}}}]{
                \texttt{\tiny{[P1$\langle$GM]}}
                \texttt{Sėkmingai perkelta 'A' į (1,1). Jūsų atnaujintas tinklelis atrodo taip:} \\
\texttt{```} \\
\texttt{ 1234567} \\
\texttt{╔═══════╗} \\
\texttt{║A◌◌◌◌ST║ 1} \\
\texttt{║P◌◌◌◌◌◌║ 2} \\
\texttt{║◌◌◌D◌┌─╢ 3} \\
\texttt{║◌◌◌┌─┤◌║ 4} \\
\texttt{║O◌◌│◌├─╢ 5} \\
\texttt{╟───┼─┘◌║ 6} \\
\texttt{║◌M◌│◌◌◌║ 7} \\
\texttt{╚═══╧═══╝} \\
\texttt{```} \\
\texttt{Šiuo metu žaidžiate 5 raundą iš maksimaliai 28.} \\
\texttt{Jūs bendrai sukaupėte 1 iš 16 baudų.} \\
\texttt{Kitas žaidėjas perkėlė objektą savo tinklelyje.} \\
\\ 
\\ 
\texttt{Kokia jūsų kita komanda?} \\
            }
        }
    }
    & \\ \\

    \theutterance \stepcounter{utterance}  
    & \multicolumn{4}{p{0.6\linewidth}}{
        \cellcolor[rgb]{0.9,0.9,0.9}{
            \makecell[{{p{\linewidth}}}]{
                \texttt{\tiny{[P1$\rangle$GM]}}
                \texttt{JUDĖTI: D, (2, 1)} \\
            }
        }
    }
    & & \\ \\

    \theutterance \stepcounter{utterance}  
    & & & \multicolumn{2}{p{0.3\linewidth}}{
        \cellcolor[rgb]{0.9,0.9,0.9}{
            \makecell[{{p{\linewidth}}}]{
                \texttt{\tiny{[GM$|$GM]}}
                \texttt{JUDĖTI: D, (2, 1)} \\
            }
        }
    }
    & & \\ \\

    \theutterance \stepcounter{utterance}  
    & & & \multicolumn{2}{p{0.3\linewidth}}{
        \cellcolor[rgb]{0.9,0.9,0.9}{
            \makecell[{{p{\linewidth}}}]{
                \texttt{\tiny{[GM$|$GM]}}
                \texttt{D moved to (2, 1)} \\
            }
        }
    }
    & & \\ \\

    \theutterance \stepcounter{utterance}  
    & & \multicolumn{4}{p{0.6\linewidth}}{
        \cellcolor[rgb]{0.9,0.9,0.9}{
            \makecell[{{p{\linewidth}}}]{
                \texttt{\tiny{[GM$\rangle$P2]}}
                \texttt{Sėkmingai perkelta 'O' į (4,1). Jūsų atnaujintas tinklelis atrodo taip:} \\
\texttt{```} \\
\texttt{ 1234567} \\
\texttt{╔═══════╗} \\
\texttt{║ADMO◌S◌║ 1} \\
\texttt{║◌◌◌◌◌◌◌║ 2} \\
\texttt{║◌◌◌◌◌┌─╢ 3} \\
\texttt{║◌◌◌┌─┤◌║ 4} \\
\texttt{║◌◌◌│◌├─╢ 5} \\
\texttt{╟───┼─┘◌║ 6} \\
\texttt{║◌◌T│P◌◌║ 7} \\
\texttt{╚═══╧═══╝} \\
\texttt{```} \\
\texttt{Šiuo metu žaidžiate 6 raundą iš maksimaliai 28.} \\
\texttt{Jūs bendrai sukaupėte 1 iš 16 baudų.} \\
\texttt{Kitas žaidėjas perkėlė objektą savo tinklelyje.} \\
\\ 
\\ 
\texttt{Kokia jūsų kita komanda?} \\
            }
        }
    }
    & \\ \\

    \theutterance \stepcounter{utterance}  
    & & & \multicolumn{4}{p{0.6\linewidth}}{
        \cellcolor[rgb]{0.9,0.9,0.9}{
            \makecell[{{p{\linewidth}}}]{
                \texttt{\tiny{[GM$\langle$P2]}}
                \texttt{JUDĖTI: P, (5, 1)} \\
            }
        }
    }
     \\ \\

    \theutterance \stepcounter{utterance}  
    & & & \multicolumn{2}{p{0.3\linewidth}}{
        \cellcolor[rgb]{0.9,0.9,0.9}{
            \makecell[{{p{\linewidth}}}]{
                \texttt{\tiny{[GM$|$GM]}}
                \texttt{JUDĖTI: P, (5, 1)} \\
            }
        }
    }
    & & \\ \\

    \theutterance \stepcounter{utterance}  
    & & & \multicolumn{2}{p{0.3\linewidth}}{
        \cellcolor[rgb]{0.9,0.9,0.9}{
            \makecell[{{p{\linewidth}}}]{
                \texttt{\tiny{[GM$|$GM]}}
                \texttt{P moved to (5, 1)} \\
            }
        }
    }
    & & \\ \\

    \theutterance \stepcounter{utterance}  
    & & \multicolumn{4}{p{0.6\linewidth}}{
        \cellcolor[rgb]{0.9,0.9,0.9}{
            \makecell[{{p{\linewidth}}}]{
                \texttt{\tiny{[P1$\langle$GM]}}
                \texttt{Sėkmingai perkelta 'D' į (2,1). Jūsų atnaujintas tinklelis atrodo taip:} \\
\texttt{```} \\
\texttt{ 1234567} \\
\texttt{╔═══════╗} \\
\texttt{║AD◌◌◌ST║ 1} \\
\texttt{║P◌◌◌◌◌◌║ 2} \\
\texttt{║◌◌◌◌◌┌─╢ 3} \\
\texttt{║◌◌◌┌─┤◌║ 4} \\
\texttt{║O◌◌│◌├─╢ 5} \\
\texttt{╟───┼─┘◌║ 6} \\
\texttt{║◌M◌│◌◌◌║ 7} \\
\texttt{╚═══╧═══╝} \\
\texttt{```} \\
\texttt{Šiuo metu žaidžiate 6 raundą iš maksimaliai 28.} \\
\texttt{Jūs bendrai sukaupėte 1 iš 16 baudų.} \\
\texttt{Kitas žaidėjas perkėlė objektą savo tinklelyje.} \\
\\ 
\\ 
\texttt{Kokia jūsų kita komanda?} \\
            }
        }
    }
    & \\ \\

    \theutterance \stepcounter{utterance}  
    & \multicolumn{4}{p{0.6\linewidth}}{
        \cellcolor[rgb]{0.9,0.9,0.9}{
            \makecell[{{p{\linewidth}}}]{
                \texttt{\tiny{[P1$\rangle$GM]}}
                \texttt{JUDĖTI: M, (3, 1)} \\
            }
        }
    }
    & & \\ \\

    \theutterance \stepcounter{utterance}  
    & & & \multicolumn{2}{p{0.3\linewidth}}{
        \cellcolor[rgb]{0.9,0.9,0.9}{
            \makecell[{{p{\linewidth}}}]{
                \texttt{\tiny{[GM$|$GM]}}
                \texttt{JUDĖTI: M, (3, 1)} \\
            }
        }
    }
    & & \\ \\

    \theutterance \stepcounter{utterance}  
    & & & \multicolumn{2}{p{0.3\linewidth}}{
        \cellcolor[rgb]{0.9,0.9,0.9}{
            \makecell[{{p{\linewidth}}}]{
                \texttt{\tiny{[GM$|$GM]}}
                \texttt{M moved to (3, 1)} \\
            }
        }
    }
    & & \\ \\

    \theutterance \stepcounter{utterance}  
    & & \multicolumn{4}{p{0.6\linewidth}}{
        \cellcolor[rgb]{0.9,0.9,0.9}{
            \makecell[{{p{\linewidth}}}]{
                \texttt{\tiny{[GM$\rangle$P2]}}
                \texttt{Sėkmingai perkelta 'P' į (5,1). Jūsų atnaujintas tinklelis atrodo taip:} \\
\texttt{```} \\
\texttt{ 1234567} \\
\texttt{╔═══════╗} \\
\texttt{║ADMOPS◌║ 1} \\
\texttt{║◌◌◌◌◌◌◌║ 2} \\
\texttt{║◌◌◌◌◌┌─╢ 3} \\
\texttt{║◌◌◌┌─┤◌║ 4} \\
\texttt{║◌◌◌│◌├─╢ 5} \\
\texttt{╟───┼─┘◌║ 6} \\
\texttt{║◌◌T│◌◌◌║ 7} \\
\texttt{╚═══╧═══╝} \\
\texttt{```} \\
\texttt{Šiuo metu žaidžiate 7 raundą iš maksimaliai 28.} \\
\texttt{Jūs bendrai sukaupėte 1 iš 16 baudų.} \\
\texttt{Kitas žaidėjas perkėlė objektą savo tinklelyje.} \\
\\ 
\\ 
\texttt{Kokia jūsų kita komanda?} \\
            }
        }
    }
    & \\ \\

    \theutterance \stepcounter{utterance}  
    & & & \multicolumn{4}{p{0.6\linewidth}}{
        \cellcolor[rgb]{0.9,0.9,0.9}{
            \makecell[{{p{\linewidth}}}]{
                \texttt{\tiny{[GM$\langle$P2]}}
                \texttt{JUDĖTI: T, (7, 1)} \\
            }
        }
    }
     \\ \\

    \theutterance \stepcounter{utterance}  
    & & & \multicolumn{2}{p{0.3\linewidth}}{
        \cellcolor[rgb]{0.9,0.9,0.9}{
            \makecell[{{p{\linewidth}}}]{
                \texttt{\tiny{[GM$|$GM]}}
                \texttt{JUDĖTI: T, (7, 1)} \\
            }
        }
    }
    & & \\ \\

    \theutterance \stepcounter{utterance}  
    & & & \multicolumn{2}{p{0.3\linewidth}}{
        \cellcolor[rgb]{0.9,0.9,0.9}{
            \makecell[{{p{\linewidth}}}]{
                \texttt{\tiny{[GM$|$GM]}}
                \texttt{T moved to (7, 1)} \\
            }
        }
    }
    & & \\ \\

    \theutterance \stepcounter{utterance}  
    & & \multicolumn{4}{p{0.6\linewidth}}{
        \cellcolor[rgb]{0.9,0.9,0.9}{
            \makecell[{{p{\linewidth}}}]{
                \texttt{\tiny{[P1$\langle$GM]}}
                \texttt{Sėkmingai perkelta 'M' į (3,1). Jūsų atnaujintas tinklelis atrodo taip:} \\
\texttt{```} \\
\texttt{ 1234567} \\
\texttt{╔═══════╗} \\
\texttt{║ADM◌◌ST║ 1} \\
\texttt{║P◌◌◌◌◌◌║ 2} \\
\texttt{║◌◌◌◌◌┌─╢ 3} \\
\texttt{║◌◌◌┌─┤◌║ 4} \\
\texttt{║O◌◌│◌├─╢ 5} \\
\texttt{╟───┼─┘◌║ 6} \\
\texttt{║◌◌◌│◌◌◌║ 7} \\
\texttt{╚═══╧═══╝} \\
\texttt{```} \\
\texttt{Šiuo metu žaidžiate 7 raundą iš maksimaliai 28.} \\
\texttt{Jūs bendrai sukaupėte 1 iš 16 baudų.} \\
\texttt{Kitas žaidėjas perkėlė objektą savo tinklelyje.} \\
\\ 
\\ 
\texttt{Kokia jūsų kita komanda?} \\
            }
        }
    }
    & \\ \\

    \theutterance \stepcounter{utterance}  
    & \multicolumn{4}{p{0.6\linewidth}}{
        \cellcolor[rgb]{0.9,0.9,0.9}{
            \makecell[{{p{\linewidth}}}]{
                \texttt{\tiny{[P1$\rangle$GM]}}
                \texttt{JUDĖTI: O, (4, 1)} \\
            }
        }
    }
    & & \\ \\

    \theutterance \stepcounter{utterance}  
    & & & \multicolumn{2}{p{0.3\linewidth}}{
        \cellcolor[rgb]{0.9,0.9,0.9}{
            \makecell[{{p{\linewidth}}}]{
                \texttt{\tiny{[GM$|$GM]}}
                \texttt{JUDĖTI: O, (4, 1)} \\
            }
        }
    }
    & & \\ \\

    \theutterance \stepcounter{utterance}  
    & & & \multicolumn{2}{p{0.3\linewidth}}{
        \cellcolor[rgb]{0.9,0.9,0.9}{
            \makecell[{{p{\linewidth}}}]{
                \texttt{\tiny{[GM$|$GM]}}
                \texttt{O moved to (4, 1)} \\
            }
        }
    }
    & & \\ \\

    \theutterance \stepcounter{utterance}  
    & & \multicolumn{4}{p{0.6\linewidth}}{
        \cellcolor[rgb]{0.9,0.9,0.9}{
            \makecell[{{p{\linewidth}}}]{
                \texttt{\tiny{[GM$\rangle$P2]}}
                \texttt{Sėkmingai perkelta 'T' į (7,1). Jūsų atnaujintas tinklelis atrodo taip:} \\
\texttt{```} \\
\texttt{ 1234567} \\
\texttt{╔═══════╗} \\
\texttt{║ADMOPST║ 1} \\
\texttt{║◌◌◌◌◌◌◌║ 2} \\
\texttt{║◌◌◌◌◌┌─╢ 3} \\
\texttt{║◌◌◌┌─┤◌║ 4} \\
\texttt{║◌◌◌│◌├─╢ 5} \\
\texttt{╟───┼─┘◌║ 6} \\
\texttt{║◌◌◌│◌◌◌║ 7} \\
\texttt{╚═══╧═══╝} \\
\texttt{```} \\
\texttt{Šiuo metu žaidžiate 8 raundą iš maksimaliai 28.} \\
\texttt{Jūs bendrai sukaupėte 1 iš 16 baudų.} \\
\texttt{Kitas žaidėjas perkėlė objektą savo tinklelyje.} \\
\\ 
\\ 
\texttt{Kokia jūsų kita komanda?} \\
            }
        }
    }
    & \\ \\

    \theutterance \stepcounter{utterance}  
    & & & \multicolumn{4}{p{0.6\linewidth}}{
        \cellcolor[rgb]{0.9,0.9,0.9}{
            \makecell[{{p{\linewidth}}}]{
                \texttt{\tiny{[GM$\langle$P2]}}
                \texttt{SAKYTI: Visi objektai sudėti į viršutinę eilutę: A(1,1), D(2,1), M(3,1), O(4,1), P(5,1), S(6,1), T(7,1). Baigta?} \\
            }
        }
    }
     \\ \\

    \theutterance \stepcounter{utterance}  
    & & & \multicolumn{2}{p{0.3\linewidth}}{
        \cellcolor[rgb]{0.9,0.9,0.9}{
            \makecell[{{p{\linewidth}}}]{
                \texttt{\tiny{[GM$|$GM]}}
                \texttt{SAKYTI: Visi objektai sudėti į viršutinę eilutę: A(1,1), D(2,1), M(3,1), O(4,1), P(5,1), S(6,1), T(7,1). Baigta?} \\
            }
        }
    }
    & & \\ \\

    \theutterance \stepcounter{utterance}  
    & & \multicolumn{4}{p{0.6\linewidth}}{
        \cellcolor[rgb]{0.9,0.9,0.9}{
            \makecell[{{p{\linewidth}}}]{
                \texttt{\tiny{[P1$\langle$GM]}}
                \texttt{Sėkmingai perkelta 'O' į (4,1). Jūsų atnaujintas tinklelis atrodo taip:} \\
\texttt{```} \\
\texttt{ 1234567} \\
\texttt{╔═══════╗} \\
\texttt{║ADMO◌ST║ 1} \\
\texttt{║P◌◌◌◌◌◌║ 2} \\
\texttt{║◌◌◌◌◌┌─╢ 3} \\
\texttt{║◌◌◌┌─┤◌║ 4} \\
\texttt{║◌◌◌│◌├─╢ 5} \\
\texttt{╟───┼─┘◌║ 6} \\
\texttt{║◌◌◌│◌◌◌║ 7} \\
\texttt{╚═══╧═══╝} \\
\texttt{```} \\
\texttt{Šiuo metu žaidžiate 8 raundą iš maksimaliai 28.} \\
\texttt{Jūs bendrai sukaupėte 1 iš 16 baudų.} \\
\texttt{Kitas žaidėjas atsiuntė šią žinutę:} \\
\\ 
\\ 
\texttt{"Visi objektai sudėti į viršutinę eilutę: A(1,1), D(2,1), M(3,1), O(4,1), P(5,1), S(6,1), T(7,1). Baigta?"} \\
\\ 
\\ 
\texttt{Kokia jūsų kita komanda?} \\
            }
        }
    }
    & \\ \\

    \theutterance \stepcounter{utterance}  
    & \multicolumn{4}{p{0.6\linewidth}}{
        \cellcolor[rgb]{0.9,0.9,0.9}{
            \makecell[{{p{\linewidth}}}]{
                \texttt{\tiny{[P1$\rangle$GM]}}
                \texttt{JUDĖTI: P, (5, 1)} \\
            }
        }
    }
    & & \\ \\

    \theutterance \stepcounter{utterance}  
    & & & \multicolumn{2}{p{0.3\linewidth}}{
        \cellcolor[rgb]{0.9,0.9,0.9}{
            \makecell[{{p{\linewidth}}}]{
                \texttt{\tiny{[GM$|$GM]}}
                \texttt{JUDĖTI: P, (5, 1)} \\
            }
        }
    }
    & & \\ \\

    \theutterance \stepcounter{utterance}  
    & & & \multicolumn{2}{p{0.3\linewidth}}{
        \cellcolor[rgb]{0.9,0.9,0.9}{
            \makecell[{{p{\linewidth}}}]{
                \texttt{\tiny{[GM$|$GM]}}
                \texttt{P moved to (5, 1)} \\
            }
        }
    }
    & & \\ \\

    \theutterance \stepcounter{utterance}  
    & & \multicolumn{4}{p{0.6\linewidth}}{
        \cellcolor[rgb]{0.9,0.9,0.9}{
            \makecell[{{p{\linewidth}}}]{
                \texttt{\tiny{[GM$\rangle$P2]}}
                \texttt{Jūsų žinutė buvo perduota kitam žaidėjui.} \\
\\ 
\\ 
\texttt{Šiuo metu žaidžiate 9 raundą iš maksimaliai 28.} \\
\texttt{Jūs bendrai sukaupėte 1 iš 16 baudų.} \\
\texttt{Kitas žaidėjas perkėlė objektą savo tinklelyje.} \\
\\ 
\\ 
\texttt{Kokia jūsų kita komanda?} \\
            }
        }
    }
    & \\ \\

    \theutterance \stepcounter{utterance}  
    & & & \multicolumn{4}{p{0.6\linewidth}}{
        \cellcolor[rgb]{0.9,0.9,0.9}{
            \makecell[{{p{\linewidth}}}]{
                \texttt{\tiny{[GM$\langle$P2]}}
                \texttt{SAKYTI: Ar baigei visus perkėlimus? Kai būsi pasiruošęs, atsakyk „baigta!".} \\
            }
        }
    }
     \\ \\

    \theutterance \stepcounter{utterance}  
    & & & \multicolumn{2}{p{0.3\linewidth}}{
        \cellcolor[rgb]{0.9,0.9,0.9}{
            \makecell[{{p{\linewidth}}}]{
                \texttt{\tiny{[GM$|$GM]}}
                \texttt{SAKYTI: Ar baigei visus perkėlimus? Kai būsi pasiruošęs, atsakyk „baigta!".} \\
            }
        }
    }
    & & \\ \\

    \theutterance \stepcounter{utterance}  
    & & \multicolumn{4}{p{0.6\linewidth}}{
        \cellcolor[rgb]{0.9,0.9,0.9}{
            \makecell[{{p{\linewidth}}}]{
                \texttt{\tiny{[P1$\langle$GM]}}
                \texttt{Sėkmingai perkelta 'P' į (5,1). Jūsų atnaujintas tinklelis atrodo taip:} \\
\texttt{```} \\
\texttt{ 1234567} \\
\texttt{╔═══════╗} \\
\texttt{║ADMOPST║ 1} \\
\texttt{║◌◌◌◌◌◌◌║ 2} \\
\texttt{║◌◌◌◌◌┌─╢ 3} \\
\texttt{║◌◌◌┌─┤◌║ 4} \\
\texttt{║◌◌◌│◌├─╢ 5} \\
\texttt{╟───┼─┘◌║ 6} \\
\texttt{║◌◌◌│◌◌◌║ 7} \\
\texttt{╚═══╧═══╝} \\
\texttt{```} \\
\texttt{Šiuo metu žaidžiate 9 raundą iš maksimaliai 28.} \\
\texttt{Jūs bendrai sukaupėte 1 iš 16 baudų.} \\
\texttt{Kitas žaidėjas atsiuntė šią žinutę:} \\
\\ 
\\ 
\texttt{"Ar baigei visus perkėlimus? Kai būsi pasiruošęs, atsakyk „baigta!"."} \\
\\ 
\\ 
\texttt{Kokia jūsų kita komanda?} \\
            }
        }
    }
    & \\ \\

    \theutterance \stepcounter{utterance}  
    & \multicolumn{4}{p{0.6\linewidth}}{
        \cellcolor[rgb]{0.9,0.9,0.9}{
            \makecell[{{p{\linewidth}}}]{
                \texttt{\tiny{[P1$\rangle$GM]}}
                \texttt{SAKYTI: baigta!} \\
            }
        }
    }
    & & \\ \\

    \theutterance \stepcounter{utterance}  
    & & & \multicolumn{2}{p{0.3\linewidth}}{
        \cellcolor[rgb]{0.9,0.9,0.9}{
            \makecell[{{p{\linewidth}}}]{
                \texttt{\tiny{[GM$|$GM]}}
                \texttt{SAKYTI: baigta!} \\
            }
        }
    }
    & & \\ \\

    \theutterance \stepcounter{utterance}  
    & & & \multicolumn{2}{p{0.3\linewidth}}{
        \cellcolor[rgb]{0.9,0.9,0.9}{
            \makecell[{{p{\linewidth}}}]{
                \texttt{\tiny{[GM$|$GM]}}
                \texttt{true} \\
            }
        }
    }
    & & \\ \\

    \theutterance \stepcounter{utterance}  
    & & & \multicolumn{2}{p{0.3\linewidth}}{
        \cellcolor[rgb]{0.9,0.9,0.9}{
            \makecell[{{p{\linewidth}}}]{
                \texttt{\tiny{[GM$|$GM]}}
                \texttt{Initial states:} \\
\texttt{```} \\
\texttt{Player 1:} \\
\texttt{ 1234567} \\
\texttt{╔═══════╗} \\
\texttt{║◌◌◌◌◌T◌║ 1} \\
\texttt{║P◌◌◌◌◌◌║ 2} \\
\texttt{║◌◌AD◌┌─╢ 3} \\
\texttt{║◌◌◌┌─┤S║ 4} \\
\texttt{║O◌◌│◌├─╢ 5} \\
\texttt{╟───┼─┘◌║ 6} \\
\texttt{║◌M◌│◌◌◌║ 7} \\
\texttt{╚═══╧═══╝} \\
\\ 
\\ 
\texttt{Player 2:} \\
\texttt{ 1234567} \\
\texttt{╔═══════╗} \\
\texttt{║◌◌◌◌◌S◌║ 1} \\
\texttt{║◌◌◌◌O◌◌║ 2} \\
\texttt{║◌◌◌M◌┌─╢ 3} \\
\texttt{║◌◌A┌─┤◌║ 4} \\
\texttt{║◌◌◌│◌├─╢ 5} \\
\texttt{╟───┼─┘◌║ 6} \\
\texttt{║D◌T│P◌◌║ 7} \\
\texttt{╚═══╧═══╝} \\
\\ 
\texttt{```} \\
            }
        }
    }
    & & \\ \\

    \theutterance \stepcounter{utterance}  
    & & & \multicolumn{2}{p{0.3\linewidth}}{
        \cellcolor[rgb]{0.9,0.9,0.9}{
            \makecell[{{p{\linewidth}}}]{
                \texttt{\tiny{[GM$|$GM]}}
                \texttt{End states:} \\
\texttt{```} \\
\texttt{Player 1:} \\
\texttt{ 1234567} \\
\texttt{╔═══════╗} \\
\texttt{║ADMOPST║ 1} \\
\texttt{║◌◌◌◌◌◌◌║ 2} \\
\texttt{║◌◌◌◌◌┌─╢ 3} \\
\texttt{║◌◌◌┌─┤◌║ 4} \\
\texttt{║◌◌◌│◌├─╢ 5} \\
\texttt{╟───┼─┘◌║ 6} \\
\texttt{║◌◌◌│◌◌◌║ 7} \\
\texttt{╚═══╧═══╝} \\
\\ 
\\ 
\texttt{Player 2:} \\
\texttt{ 1234567} \\
\texttt{╔═══════╗} \\
\texttt{║ADMOPST║ 1} \\
\texttt{║◌◌◌◌◌◌◌║ 2} \\
\texttt{║◌◌◌◌◌┌─╢ 3} \\
\texttt{║◌◌◌┌─┤◌║ 4} \\
\texttt{║◌◌◌│◌├─╢ 5} \\
\texttt{╟───┼─┘◌║ 6} \\
\texttt{║◌◌◌│◌◌◌║ 7} \\
\texttt{╚═══╧═══╝} \\
\\ 
\texttt{```} \\
            }
        }
    }
    & & \\ \\

    \theutterance \stepcounter{utterance}  
    & & & \multicolumn{2}{p{0.3\linewidth}}{
        \cellcolor[rgb]{0.9,0.9,0.9}{
            \makecell[{{p{\linewidth}}}]{
                \texttt{\tiny{[GM$|$GM]}}
                \texttt{* success: True} \\
\texttt{* lose: False} \\
\texttt{* aborted: False} \\
\texttt{{-}{-}{-}{-}{-}{-}{-}} \\
\texttt{* Shifts: 12.00} \\
\texttt{* Max Shifts: 12.00} \\
\texttt{* Min Shifts: 6.00} \\
\texttt{* End Distance Sum: 0.00} \\
\texttt{* Init Distance Sum: 31.75} \\
\texttt{* Expected Distance Sum: 29.33} \\
\texttt{* Penalties: 1.00} \\
\texttt{* Max Penalties: 16.00} \\
\texttt{* Rounds: 9.00} \\
\texttt{* Max Rounds: 28.00} \\
\texttt{* Object Count: 7.00} \\
            }
        }
    }
    & & \\ \\

\end{supertabular}
}

\end{document}
